\subsection{Allocating the full positive mean jointly with the comparison}
\label{subsec:joint-mean}

The closing inequality contains the comparison energy with coefficient
$63/64$ and the mean energy with coefficient $-1/64$. A positive
diagonal occurring in both can therefore reduce the closing bound:
subtract its mean coefficient from $63$ times its comparison coefficient,
then apply an upper bound to the remaining nonnegative energy.
This is the joint allocation of \cref{cor:joint-closing}.
We apply it to the positive expansion in Step~2 of
\cref{prop:record:mean-credit}. Of the lower bounds retained there
(called mean credits), only those supported on one block are used
separately here. The following argument replaces the credits on
two or more blocks.

For a middle block write
\[
 \alpha_B=\frac{w_{{\rm D},B}r_B^2}{1+h_B},\qquad
 T_B=E_{{\rm add},B}+H_B
     =\beta_B(D_{E_0,B}+D_{E_1,B}),\qquad
 \mathsf E_B=E_{{\rm add},B}+G_B.
\]
The measures $E_{{\rm add},B},G_B,H_B$, the fold diagonals $D_{F_j,B}$,
and the matching diagonals $D_{E_j,B}$ are those defined in
\cref{subsec:record:positive_means}; $K_{s,B}$ is the added diagonal
on side $s\in\{{\rm L},{\rm R}\}$ of the matching bracket
(\cref{prop:record:matching}). Here $\beta_B$ is a local coefficient.
The actual chosen diagonals satisfy
\begin{equation}\label{eq:joint-resources}
 \begin{aligned}
 2\le\beta_B/\alpha_B&\le R=2.004317,\qquad\sigma=0.3842,\\
 Q_{E,B}&\succeq_\iota\alpha_B(D_{E_0,B}+D_{E_1,B}),\\
 K_{s,B}&\succeq_\iota\sigma\alpha_B(D_{F_1,B}+D_{F_2,B}).
 \end{aligned}
\end{equation}
The density of an actual fold marginal relative to its normalized
companion is at most $H_*=1.000003310$. Thus its companion diagonal
dominates $H_*^{-1}D_{F_j,B}$.
These assertions use the same $33$ physical pieces and chosen sides
as \cref{prop:record:matching}. Explicitly, on a piece starting at $X$,
\[
 \frac{\beta_B}{\alpha_B}
 \le\frac{v_{\rm D}}{w_{\rm D}}
       +\frac{6}{\lfloor\sqrt X\rfloor(1-3/X)}\le R,
 \qquad
 \frac{\min(k_{{\rm L},i},k_{{\rm R},i})}{w_{{\rm D},i}H_i}\ge\sigma.
\]
The matching coefficient $q_i\ge w_{{\rm D},i}/(1+h_i)$ and
$H_i\le H_*$ complete the checks. Both independent assemblies verify
these inequalities piece by piece.

\begin{lemma}[Multiaffine allocation constants]\label{lem:joint-multiaffine}
Put $l=1/R$, $u=1/2$, and, for $0\le k\le n$, put
\[
 P_{n,k}=(1+l)^k(1+u)^{n-k},\quad
 E_{1,n,k}=kl+(n-k)u,\quad
 E_{2,n,k}=\binom k2l^2+k(n-k)lu+\binom{n-k}2u^2.
\]
Define
\[
 \eta_n=\min_{0\le k\le n}\frac{P_{n,k}-1}{E_{2,n,k}} (2\le n\le10),
 \qquad
 \zeta_n=\min_{0\le k\le n}\frac{P_{n,k}}{E_{1,n,k}} (1\le n\le9).
\]
Write $e_1(z)=\sum_i z_i$ and $e_2(z)=\sum_{i<j}z_iz_j$.
For $z_i\in[l,u]$,
$\prod_i(1+z_i)-1\ge\eta_n e_2(z)$ and
$\prod_i(1+z_i)\ge\zeta_n e_1(z)$ in the stated degree ranges.
In every degree the corresponding constants
$\eta_0=211/80$ and $\zeta_0=9/4$ are valid.
\end{lemma}

\begin{proof}
For a fixed trial constant each residual is multiaffine, so its minimum
on the box is attained at a vertex. The displayed formulas are precisely
its vertex tests, grouped by the number of lower coordinates.
For the uniform bounds enlarge the box to $[0,1/2]^n$.
At a vertex with $k$ nonzero coordinates the ratios are
$4((3/2)^k-1)/\binom k2$ and $2(3/2)^k/k$.
The first is minimized at $k=5$: its adjacent difference has the sign
of $(k-5)(3/2)^k+4$. The second is minimized at $k=2,3$,
as its adjacent difference has the sign of $k-2$.
The minima are $211/80$ and $9/4$.
\end{proof}

Use the normalized price assignments of
\cref{lem:four-routes-normalized}, repriced to the selected exponent $a$.
Let $x=A^{\rm arm}_F=A^{\rm arm}_E$ be the common raw-arm total of
\cref{def:closing-quantities}(3): the sum of the arm coefficients
before multiplying by their leading operator norms. Set $y=\sigma x$.
For this subsection abbreviate
$U=tU_0^{\prime\rm new}$ and $V=t{\Uinv}^{\rm new}$, where
$t=e^{2(a-a_0)570}$. These are the total first-stage price and the
portion assigned to invariant factors, respectively; $0\le V\le U$.
An outside factor means a factor on a block other than the selected
fold or matching blocks. Define
\begin{align*}
 F_\eta&=\eta_0e^U+\sum_{j=0}^8(\eta_{j+2}-\eta_0)V^j/j!,\\
 F_\zeta&=\zeta_0e^U+\sum_{j=0}^8(\zeta_{j+1}-\zeta_0)V^j/j!,\\
 B_\rho&=e^U+\sum_{j=1}^{10}(\rho^j-1)V^j/j!,\qquad \rho=1+1/R,\\
 w&=\frac{U}{1+\rho_H},\qquad
 \rho_H=\frac{3}{\lfloor\sqrt{Q_3}\rfloor(1-3/Q_3)}.
\end{align*}

\begin{proposition}[Four disjoint middle allocations]\label{prop:joint-middle-mean}
In the joint closing of \cref{cor:joint-closing}, the middle mean
decreases the closing number, in addition to its one-block mean
credits, by $\Delta_\alpha+\Delta_M+\Delta_O+\Delta_K$, where
\begin{align}
 \Delta_\alpha&=\frac{x^2}{16H_*^2}
                 \frac{e^w-1-w}{w^2},\label{eq:joint-alpha}\\
 \Delta_M&=\frac{x^2}{128}F_\eta,\qquad
 \Delta_O=\frac{x}{64}\left(y+\frac{y^3}{6}+\frac{y^5}{120}\right)F_\zeta,
 \qquad
 \Delta_K=\frac1{64}\left(\frac{y^2}{2}+\frac{y^4}{24}+\frac{y^6}{720}\right)B_\rho.
 \label{eq:joint-middle}
\end{align}
These are corrections to the combined comparison and mean bound;
they are not separate lower bounds obtained by substituting upper
prices in a mean energy.
\end{proposition}

\begin{proof}
After paying $R_{\rm err}e^{T_K+T_{QE}}$, the available positive product is
\[
 \bigotimes_B(\mu_B+E_{{\rm add},B}+G_B+H_B+K_{s,B}+Q_{E,B}).
\]
A tag specifies the summand chosen at a block in this tensor product.
Partition the retained terms first by their exact sets of $K_s$ tags
and then by their $Q_E$ tags. The four families are:
\begin{enumerate}[label=(\roman*)]
\item no $K_s$ or $Q_E$; at least two $\mathsf E$ tags, with reference
factors elsewhere;
\item no $K_s$; a nonempty $Q_E$ set and at least two positions carrying
$T$ or $Q_E$, with $\mu$ or $G$ at the remaining positions;
\item exactly $1,3,5$ $K_s$ tags; at least one further position carrying
$T$ or $Q_E$, with $\mu$ or $G$ at the remaining positions;
\item exactly $2,4,6$ $K_s$ tags, with arbitrary $\mu,T,Q_E,G$ outside.
\end{enumerate}
They are disjoint and correspond, in order, to
$\Delta_\alpha,\Delta_M,\Delta_O,\Delta_K$.
Each has support on at least two blocks, so none consumes a one-block
mean credit.

For the first family use
$\mathsf E_B\pm F_{{\rm fold},B}\succeq_{\vartheta_0}0$ on a selected pair.
The other $\mathsf E$ factors and reference factors are positive at
the same cut. The two companions dominate
$2\alpha_B(D_{F_1,B}+D_{F_2,B})/H_*$ after returning to the reserved
diagonal. Average the pair selection over the $\binom n2$ pairs of
an $n$-tag term. With raw coordinate $x_B=2\alpha_B$ this gives
coefficient $4/(H_*^2\binom n2)$ on the selected pair.
Consequently its exponential completion, after division by $64$, is
$x^2\sum_{j\ge0}w^j/(16H_*^2(j+2)!)$.

Here the $\mathsf E$ and $H$ assignments must be made together.
The physical bound $H_B\le\rho_H\mathsf E_B$ in positive mass allows
their total $U$ to be filled with an $\mathsf E$ assignment at least
$U/(1+\rho_H)$: first assign the genuine $H$ mass and put the remaining
slack into $\mathsf E$. The invariant and generic parts can be filled
with their original separate nonnegative slacks. This preserves $U$
and an invariant assignment $V_{\rm full}\ge V$. The completed
comparison cost is nonincreasing in $V_{\rm full}$, while the mean
corrections are nondecreasing in it, so replacing it by $V$ is safe.
No lower bound for the physical $\mathsf E$ mass is inferred from $U$.

For the second family factor out the $T$ coefficient at each matching
position and put $z_B=\alpha_B/\beta_B\in[1/R,1/2]$.
The nonempty-$Q_E$ sum dominates
$\prod T_B(\prod(1+z_B)-1)$.
The first inequality of \cref{lem:joint-multiaffine} allocates it
simultaneously over selected pairs; the factor $1/2$ in the unordered
pair completion gives $\Delta_M$.
For odd $K_s$ degree select one matching position instead and use the
second inequality. For even $K_s$ degree keep the whole outside
expansion and use $T_B+Q_{E,B}\succeq_\iota\rho T_B$ on a pure
invariant background. These give $\Delta_O$ and $\Delta_K$.
The uniform constants apply also to backgrounds containing $G$;
the finite improvements are used only on pure invariant backgrounds.
The multiaffine inequalities allocate a term among its eligible
positions; they never spend it separately in full on every position.

It remains to justify completion of the combined expression. For a
labelled diagonal energy $Z$, combine its comparison coefficient $c$
and retained mean coefficient $b$ as $(63c-b)Z$ first. The finite
checks below show $63c-b\ge0$, including the already retained variance
credits. Thus \cref{lem:record:record-expansion} applies to the residual
positive series. In particular, allowing repeated block labels enlarges
the sum over distinct blocks and gives the displayed exponential
series; this step is legitimate because the residual coefficients
are nonnegative. Its value is precisely
the former comparison bound minus the four displayed corrections.
This proves the proposition without asserting a lower bound for any
isolated upper-priced mean term.
\end{proof}

\paragraph{Residual coefficients.}
The sixteen classes and their invariant norm caps $\kappa_b$ are
those of \cref{cor:envelope-costs}. In the table below $F$ denotes a
selected fold arm and $E$ a selected matching arm. The outside degree
$j$ counts additional invariant factors. Check $j=0,\ldots,10$
separately, and use one further category for $j\ge11$.
The new mean coefficient is the sum of the applicable entries
\[
\begin{array}{c|c}
\text{selected class}&\text{mean coefficient}\\ \hline
F^2&4/H_*^2\\
E^2&\eta_{j+2}\ (j\le8),\quad\eta_0\ \text{otherwise}\\
F^nE,\ n=1,3,5&\sigma^n\zeta_{j+1}\ (j\le8),\quad\sigma^n\zeta_0\ \text{otherwise}\\
F^n,\ n=2,4,6&\sigma^n\rho^j\ (j\le10),\quad\sigma^n\ \text{in the tail}.
\end{array}
\]
The $F^2$ entries are added. Every selected block cap is at least $2$.
The restoration formulas therefore bound a selected-pair gain by $1/4$;
the same bound holds coefficientwise on each matching background.
Subtract at most $\kappa_b/4$ for $F^2$ and aligned $E^2$,
$\kappa_b/32$ for $F^4$, and the aligned $E^4$ cap divided by $32$.
All $16\cdot12=192$ residuals $63(\kappa_b-\text{variance})-
\text{mean}$ are positive; the minimum exceeds $23.66$.
The gates also hold before averaging over the record words. For the pure
fold classes $F^n$, $n=2,4,6$, the cap $\kappa_n$ is common to every word:
the chart reflection is independent of the word, and the extra reflection
acts trivially when there is no matching coordinate. In addition to the
class gates, check
\[
 63\bigl(1-\tfrac14\one_{n\in\{2,4\}}\bigr)\kappa_n
 -\frac4{H_*^2}\one_{n=2}-\sigma^n\rho^j>0
 \qquad(0\le j\le10),
\]
and the same inequality with $\rho^j$ replaced by $1$ for the tail.
These $36$ gates allow the full conservative variance credit on an
aligned quartic word, before its weight $1/8$ in the class is taken.
Each matching-pair class has a common cap; the pure matching quartics
have no new mean subtraction and retain their existing variance gates.
For the mixed words the chosen unframed caps are at least $1/4$,
and $63/4-\sigma^n\zeta_{j+1}>0$ for $n=1,3,5$, $0\le j\le8$.
The exact rational lists and their independent reconstruction are
part of \cref{cert:four-routes-closing}.

\begin{proposition}[Mean terms with a nonempty upper support]\label{prop:joint-upper-mean}
Let $S_u,A_u$ denote just the finite upper comparison assignments,
excluding prefix and tail padding. Set
\[
 C_E=e^{S_u}\cosh A_u-1,\qquad C_O=e^{S_u}\sinh A_u.
\]
The further joint correction is $\Delta_u-\mathcal R_u$, where
\begin{align*}
 \Delta_u={}&\frac1{64}\left\{
 C_E\bigl(\cosh(y)B_\rho+x\sinh(y)F_\zeta\bigr)
 +C_O\bigl(\sinh(y)B_\rho+x\cosh(y)F_\zeta\bigr)\right\},\\
 \mathcal R_u={}&\frac{R_{\rm err}e^{T_K+T_{QE}}}{64}
                   (e^{S_u+A_u}-1).
\end{align*}
\end{proposition}

\begin{proof}
On each finite upper block both matching brackets are $\mu_B+Q_B$,
with the same positive reserved-diagonal norm envelope $Q_B$.
The nonempty-$Q$ part of their tensor product is disjoint from all
middle-only terms used above. Fix its skew parity and the exact middle
$K_s$ set. If the parities agree, retain the full $T/Q_E/G$ background;
otherwise select one additional matching position. The same multiaffine
allocation gives the displayed even and odd exponential sums.
The new first-stage error is bounded by $\mathcal R_u$, obtained by
summing over the entire nonempty upper support, including discarded
branches. Every such term has positive upper degree, so it belongs to
none of the sixteen exceptional middle classes or their variance credits.
After splitting each full middle skew price into its raw arm and
positive remainder, the residual gates are
$63-\sigma^n\rho^j>0$ and $63-\sigma^n\zeta_{j+1}>0$, with the same
finite ranges and uniform tails. Since $\sigma<1$, $n=0$ is the worst
case, and the stated rational constants satisfy these gates.
Complete this positive residual series as above. The finite upper
assignments are used jointly in its positive cost and mean term;
they are not lower bounds for the mass of $Q_B$.
\end{proof}
